Similar to the ADF scheme, adaptive sampling also tries to utilise the information gathered throughout the MD sampling process  to guide the subsequent MD moves, but differently, it applies data analytic tools such as dimensionality reductions, clustering analysis and Markov state model to identify poorly sampled configuration spaces, towards which the subsequent samplings will be biased. For instance, by building up  Markov state model  from  the collected MD frames and analysing the transition matrix or using cluster analysis, \cite{zimmerman2015fast,doerr2016htmd} one could identify microstates that have not been frequently visited in the MD runs, which can be used as the starting point for new MD simulations to enhance the exploitation on a particular part of the energy landscape. For large biomolecular simulations, a contact matrix\cite{weber2011characterization,doerr2014fly} that counts the number of times states $i$ and $j$ were adjacent to each others in  the trajectory can be built to identify states that have the minimum contacts, which are selected as new starting points for MD simulations. The idea of structural-based adaptive sampling has important implication that underscores, for example, the development of on-the-fly machine learning forcefields\cite{jinnouchi2019fly}  to accelerate \emph{ab initio} MD simulations, which uses an agnostic structural descriptor and similarity kernel\cite{bartok2013representing} to identify representative structures on the PES for forcefield training. As the final example, \citet{perez2020adaptivebandit} developed the AdaptiveBandit framework for biomolecular simulations to tackle the exploitation-exploration dilemma, in which the decision to take an action $a$ to restart the sampling from the starting point at a given time $\bm{R}(t)$ is cast into a maximisation problem over all possible actions $\mathcal{A}$:
\begin{equation*}
    a(t)=\underset{a\in\mathcal{A}}{\mbox{argmax}} \left[\mathcal{Q}_t(a)+c\sqrt{\frac{\ln t}{N_{t}(a)}}\right],
\end{equation*}
in which $\mathcal{Q}_t(a)$ is the estimated reward for taking the action $a$, $e.g.$ the likelihood that a new structure will be generated, $c$ is a parameter aim to balance  exploitation-exploration, and $N_t (a)$ is the number of time that the action $a$ has been taken up to time $t$.
